\documentclass[a4paper]{article}

% Basic packages
% Platform-aware CJK font selection: detects PingFang.ttc for macOS (fontset=mac),
% falls back to fandol (free CJK). Use PassOptionsToPackage because \IfFileExists
% cannot be embedded inside \usepackage options (would cause "Missing \endcsname").
\IfFileExists{/System/Library/Fonts/PingFang.ttc}{
    \PassOptionsToPackage{fontset=mac}{ctex}
}{
    \PassOptionsToPackage{fontset=fandol}{ctex}
}
\usepackage{ctex}
% CJK fallback for rare glyphs (e.g. 珺/昇/騰) absent from the default font.
% On macOS, Heiti SC covers nearly all CJK characters including rare ones.
% On Linux, install Noto Sans CJK and change the font name below.
\IfFileExists{/System/Library/Fonts/STHeiti Medium.ttc}{
    \newCJKfontfamily\fallbackhei{Heiti SC}
}{}
\usepackage{amsmath, amssymb}
\usepackage{graphicx}
\usepackage[margin=2.5cm]{geometry}
\usepackage[most]{tcolorbox}
\usepackage{etoolbox}
\usepackage{listings}
\usepackage{hyperref}
\usepackage{booktabs}
\usepackage{subcaption}
\usepackage{float}
\usepackage{tikz}
\usetikzlibrary{positioning}
\IfFileExists{pgfplots.sty}{
    \usepackage{pgfplots}
    \pgfplotsset{compat=1.18}
}{}

% Highlight boxes
\newtcolorbox{knowledgebox}[1]{
    enhanced, colback=blue!5!white, colframe=blue!75!black, colbacktitle=blue!75!black,
    coltitle=white, fonttitle=\bfseries, title=#1, attach boxed title to top left={yshift=-2mm, xshift=2mm},
    boxrule=1pt, sharp corners
}

\newtcolorbox{importantbox}[1]{
    enhanced, colback=yellow!10!white, colframe=yellow!80!black, colbacktitle=yellow!80!black,
    coltitle=black, fonttitle=\bfseries, title=#1, sharp corners
}

\newtcolorbox{warningbox}[1]{
    enhanced, colback=red!5!white, colframe=red!75!black, colbacktitle=red!75!black,
    coltitle=white, fonttitle=\bfseries, title=#1, sharp corners
}

% Listings
\lstset{
    language=Python,
    basicstyle=\ttfamily\small,
    keywordstyle=\color{blue},
    stringstyle=\color{red!60!black},
    commentstyle=\color{green!60!black},
    breaklines=true,
    frame=single,
    numbers=left,
    numberstyle=\tiny\color{gray},
    captionpos=b,
    extendedchars=true,
    inputencoding=utf8
}

% Front-page metadata
\newcommand{\notetitle}{课程笔记：[在此填写课程主题]}
\newcommand{\noteauthors}{五道口纳什 \& Codex}
\newcommand{\notedate}{\today}
\newcommand{\videochannel}{[在此填写频道或作者]}
\newcommand{\videopublishdate}{[在此填写发布日期]}
\newcommand{\videoduration}{[在此填写视频时长]}
\newcommand{\videourl}{}
\newcommand{\videocoverpath}{}

\begin{document}

\begin{titlepage}
\centering
{\Large 课程笔记\par}
\vspace{1.2cm}
{\huge\bfseries \notetitle\par}
\vspace{0.8cm}
{\large \noteauthors\par}
\vspace{0.3cm}
{\large \notedate\par}
\vspace{1.2cm}

\ifx\videocoverpath\empty
{\small 请在生成笔记时填入视频封面图的本地路径。\par}
\else
\includegraphics[width=0.82\textwidth,height=0.45\textheight,keepaspectratio]{\videocoverpath}\par
\fi

\vfill
\begin{tcolorbox}[width=0.9\textwidth, colback=black!2!white, colframe=black!60, sharp corners]
\textbf{视频作者/频道}：\videochannel\par
\textbf{发布日期}：\videopublishdate\par
\textbf{视频时长}：\videoduration\par
\textbf{视频链接}：
\ifx\videourl\empty
未填写
\else
\href{\videourl}{\nolinkurl{\videourl}}
\fi
\end{tcolorbox}
\end{titlepage}

\tableofcontents
\newpage

%% --- 正文内容开始 --- %%

% Replace this block with the generated notes.
% Fill the metadata block above, especially \videocoverpath, so the front page
% contains the original video cover image.
% Keep images outside knowledgebox/importantbox/warningbox.
% For any figure that comes from a video frame or a crop of a frame,
% include the source time interval as a bottom footnote on the same page.
% A stable pattern is:
% \begin{figure}[H]
% \centering
% \includegraphics[width=\textwidth]{...}
% \caption{... \protect\footnotemark}
% \end{figure}
% \footnotetext{视频画面时间区间：00:12:31--00:12:46。若为裁剪图，时间区间仍对应原视频画面。}
% End the document with a final section such as:
% \section{总结与延伸}
% covering the speaker's closing discussion and your own synthesis.
%
% --- RARE CJK GLYPHS ---
% If a character (e.g. 珺 in 张小珺, 昇/騰 in AI chip names) renders as □,
% wrap it in \fallbackhei{}:  {\fallbackhei 张小珺}
% After compilation, always check for missing glyphs:
%   grep -c "Missing character" <basename>.log
% Non-zero = some characters didn't render; find them with:
%   grep "Missing character" <basename>.log | sort -u

%% --- 正文内容结束 --- %%

\end{document}
